26 2. Hilbert spaces
Note that (almost) the same proof shows that $K$ is compact when defined
on $\mathcal{C}[a, b]$.
\begin{theorem}{2.15 (Arzelà-Ascoli).}
Suppose the sequence of functions $f_n(x)$,
$n\in \mathbb{N}$, on a compact interval is (uniformly) equicontinuous, that is, for
every $\varepsilon > 0$ there is a $\delta > 0$ (independent of $n$) such that
$$|f_n(x) - f_n(y)| \leq \varepsilon \quad \text{if } |x - y| < \delta.$$
(2.40)
If the sequence $f_n$ is bounded, then there is a uniformly convergent subse-
quence.
\end{theorem}
(The proof is not difficult but I still don't want to repeat it here since it
is covered in most real analysis courses.)
Compact operators are very similar to (finite) matrices as we will see in
the next section.
\begin{problem}{}
Show that compact operators form an ideal.
\end{problem}
\subsection{The spectral theorem for compact
symmetric operators}
Let $\tilde{H}$ be an inner product space. A linear operator $A$ is called symmetric
if its domain is dense and if
$$\langle g, Af \rangle = \langle Ag, f \rangle \quad f, g \in \mathcal{D}(A).$$
(2.41)
A number $z \in \mathbb{C}$ is called eigenvalue of $A$ if there is a nonzero vector
$u \in \mathcal{D}(A)$ such that
$$Au = zu.$$
(2.42)
The vector $u$ is called a corresponding eigenvector in this case. An eigen-
value is called simple if there is only one linearly independent eigenvector.
\begin{theorem}{2.16.}
Let $A$ be symmetric. Then all eigenvalues are real and
eigenvectors corresponding to different eigenvalues are orthogonal.
\end{theorem}
\begin{proof}
Suppose $\lambda$ is an eigenvalue with corresponding normalized eigen-
vector $u$. Then $\lambda = \langle u, Au \rangle = \langle Au, u \rangle = \overline{\lambda}$, which shows that $\lambda$ is real.
Furthermore, if $A u_1 = \lambda_1 u_1, j = 1, 2$, we have
$$(\lambda_1 - \lambda_2)\langle u_1, u_2 \rangle = \langle A u_1, u_2 \rangle - \langle u_1, A u_2 \rangle = 0$$
(2.43)
finishing the proof.
\end{proof}
Now we show that $A$ has an eigenvalue at all (which is not clear in the
infinite dimensional case)!
\begin{theorem}{2.17.}
A symmetric compact operator has an eigenvalue $\alpha_1$ which
satisfies $|\alpha_1| = \|A\|$.
\end{theorem}